% =============================================================
%  Appendix A5 -- Expectation Scaling Sensitivity
%  Label: app:scaling
% =============================================================

\section{Expectation Scaling Sensitivity}\label{app:scaling}

\subsection{The Scaling Problem}

The CFPS inflation expectation question is ordinal.
Respondents report whether they expect prices to rise, stay
the same, or fall.
The raw response is coded as $\mu_{it} \in \{-1, 0, 1\}$.
To construct the forecast error
$\mathrm{FE}_{ipt} = \pi^{\text{headline}}_{p,\,t\to t+1}
- \tilde{\mu}_{ipt}$,
I must convert the ordinal expectation into units comparable
with the realised CPI (which is measured in annual
percentage points).
This conversion requires a scaling function
$\tilde{\mu} = f(\mu)$ that maps ordinal codes to CPI units.

The scaling is inherently underdetermined.
The CFPS question asks for the direction of expected price
changes, not for a point forecast.
A respondent who reports ``prices will rise'' could mean
1 percent, 5 percent, or 20 percent.
No information in the survey resolves this ambiguity.
The scaling choice therefore affects the level and variance
of the forecast error and, consequently, the magnitude
(though not the sign) of the regression coefficient
$\hat{\beta}_3$.

This section evaluates the sensitivity of the main results
to the scaling choice by considering four alternative
mappings and providing economic reasoning for each.

\subsection{Scaling Alternatives}

I consider the following four scaling functions.

Baseline: $\{-2, 0, 2\}$.
The baseline maps ``fall'' to $-2$ percentage points,
``same'' to $0$, and ``rise'' to $+2$ percentage points.
The rationale is that unconditional headline CPI inflation
in China during the sample period (2016--2022) averaged
approximately 2 percent per year.
A respondent who expects ``prices will rise'' is interpreted
as expecting roughly average inflation.
A respondent who expects ``prices will fall'' is interpreted
as expecting deflation of comparable magnitude.
The symmetry of the mapping reflects the symmetric phrasing
of the survey question.

Alternative 1: $\{-1, 0, 1\}$ (unit ordinal).
The simplest scaling uses the ordinal codes directly.
The forecast error becomes realised CPI in percentage points
minus the ordinal code ($-1$, $0$, or $1$).
This scaling implicitly assumes that ``prices will rise''
corresponds to 1 percent expected inflation.
Given that average realised inflation is approximately 2
percent, this mapping understates the expected inflation of
optimistic respondents and compresses the forecast-error
distribution.

Under this scaling, the dependent variable has a narrower
range.
The forecast error for a respondent in a province with 2
percent realised inflation who reported ``rise'' is
$2 - 1 = 1$ percentage point, compared with $2 - 2 = 0$
under the baseline.
The coefficient $\hat{\beta}_3$ will be larger in magnitude
under the unit scaling because the dependent variable
attributes less of the realised CPI to the expectation
component, making the forecast error larger on average for
``rise'' respondents.

Alternative 2: $\{-4, 0, 4\}$ (wide scaling).
This scaling assumes that ``prices will rise'' implies
roughly 4 percent expected inflation.
The rationale is that during high-inflation episodes
(such as the 2019--2020 food-price spike), households
reporting ``rise'' may have had in mind inflation rates
well above the unconditional mean.
If respondents anchor on recent salient price increases
rather than on average CPI, a wider scaling may better
approximate the implied expectation.

Under this scaling, the forecast error for a ``rise''
respondent in a 2-percent-inflation province is
$2 - 4 = -2$ percentage points.
The wider scaling mechanically increases the fraction of
negative forecast errors among ``rise'' respondents,
reducing the variance of the forecast error conditional
on the shock and potentially attenuating $\hat{\beta}_3$.

Alternative 3: $\{-1, 0, 3\}$ (asymmetric).
This scaling allows the ``rise'' response to carry a
different magnitude than the ``fall'' response.
The rationale is that China experienced positive average
inflation during the sample period, so ``prices will rise''
may be the default expectation (consistent with the fact
that approximately 60 to 70 percent of respondents choose
this category in most waves).
If ``rise'' is the typical response, the associated
expected inflation may be higher than the expected
deflation associated with ``fall.''
The asymmetric mapping assigns $+3$ to ``rise'' and $-1$ to
``fall,'' reflecting the possibility that respondents who
expect prices to rise anticipate moderate inflation while
those who expect prices to fall anticipate only mild
deflation.

\subsection{Results Under Alternative Scalings}

Table~\ref{tab:scaling_sensitivity} reports the
forecast-error coefficient $\hat{\beta}_3$ under each scaling
for the preferred region-by-wave FE specification.

\begin{table}[htbp]
\centering
\caption{Forecast-Error Coefficient Under Alternative
  Expectation Scalings}
\label{tab:scaling_sensitivity}
\begin{threeparttable}
\begin{tabular}{lccccc}
\toprule
Scaling & $\tilde{\mu}$ map &
  $\hat{\beta}_3$ & SE & $p$-value & $N$ \\
\midrule
Baseline       & $\{-2,\;0,\;2\}$  & $-6.067$ & 1.145
  & $<0.001$ & 69,533 \\
Unit ordinal   & $\{-1,\;0,\;1\}$  & $-6.309$ & 1.319
  & $<0.001$ & 69,533 \\
Wide           & $\{-4,\;0,\;4\}$  & $-5.583$ & 0.978
  & $<0.001$ & 69,533 \\
Asymmetric     & $\{-1,\;0,\;3\}$  & $-5.825$ & 1.067
  & $<0.001$ & 69,533 \\
\bottomrule
\end{tabular}
\begin{tablenotes}[flushleft]
\footnotesize
\item Region-by-wave fixed effects. Standard errors clustered
  at the province level.
\item All specifications include the same set of controls as
  the preferred specification in
  Table~\ref{tab:eq3_forecast_error}.
\end{tablenotes}
\end{threeparttable}
\end{table}

The sign of $\hat{\beta}_3$ is negative and statistically
significant at $p < 0.001$ under all four scalings.
The point estimate ranges from $-5.583$ (wide scaling) to
$-6.309$ (unit ordinal scaling), a spread of approximately
12 percent around the baseline.
The variation in magnitude follows the logic outlined above.
A wider positive scaling (attributing more expected inflation
to ``rise'' respondents) mechanically reduces the forecast
error and compresses $\hat{\beta}_3$ toward zero, while a
narrower scaling inflates the forecast error and increases
$|\hat{\beta}_3|$.

\subsection{Why the Sign Is Invariant}

The invariance of the sign to scaling has a simple
algebraic explanation.
The forecast error is
$\mathrm{FE}_{ipt} = \pi_{p,\,t\to t+1} - f(\mu_{it})$.
The meat shock enters the regression through its covariance
with the forecast error.
The realised CPI $\pi_{p,\,t\to t+1}$ does not depend on the
scaling (it is an external variable), so the only channel
through which scaling affects $\hat{\beta}_3$ is through
its effect on $f(\mu_{it})$.

Since the ordinal expectation $\mu_{it}$ is positively
correlated with the meat shock (households in high-shock
provinces are more likely to report ``rise''), the scaled
expectation $f(\mu_{it})$ is also positively correlated
with the shock for any scaling that is monotonically increasing
in $\mu$ (i.e., any scaling where ``rise'' maps to a higher
value than ``fall'').
The forecast error $\pi - f(\mu)$ therefore has a negative
partial correlation with the shock because the expectation
component overshoots the CPI component.
This negative partial correlation produces a negative
$\hat{\beta}_3$ regardless of the specific values assigned
to each ordinal category, as long as the ordering is preserved.

The magnitude of $\hat{\beta}_3$ changes with the scaling
because the scaling affects the variance of $f(\mu)$ and its
covariance with the shock.
A wider scaling increases the variance of the expectation
component, which reduces the forecast error for ``rise''
respondents (making it less negative) and increases it for
``fall'' respondents (making it more positive).
The net effect on $\hat{\beta}_3$ depends on the relative
magnitudes of these adjustments, but the sign is determined
by the ordering of the categories, which is invariant to
monotone transformations.

\subsection{Which Scaling Is Most Defensible?}

No scaling can be verified from the data because the survey
does not elicit a point forecast.
However, three considerations favour the baseline
$\{-2, 0, 2\}$ mapping.

First, the baseline is centred on the unconditional mean of
realised headline CPI during the sample period
(approximately 2 percent per year).
This makes the forecast error approximately mean-zero
for respondents who correctly predict the direction of
inflation, which is a desirable property for interpreting
the constant in the regression.

Second, the symmetric mapping avoids imposing an asymmetry
on the expectation distribution that the survey question does
not support.
The CFPS question is phrased symmetrically (``will prices
rise, stay the same, or fall?''), so there is no a priori
reason to assign different magnitudes to the ``rise'' and
``fall'' categories.
The asymmetric scaling $\{-1, 0, 3\}$ is motivated by the
empirical distribution of responses (most respondents choose
``rise''), but this distribution could reflect either the
level of expected inflation or a cognitive default toward
the modal response.
Without additional information, the symmetric mapping is
more conservative.

Third, the baseline scaling produces a coefficient magnitude
($-6.067$) that is intermediate in the range across scalings.
This means that the baseline is neither the most nor the least
conservative interpretation of the overreaction.
The qualitative comparison with the rational benchmark
($|\beta_3^R| \leq 0.05$) is unaffected: the estimated
coefficient exceeds the bound by a factor of at least 112
(under the wide scaling) and at most 126 (under the unit
ordinal scaling).
The two-orders-of-magnitude excess is invariant to any
reasonable scaling choice.

\subsection{Implications for the Implied Diagnostic Parameter}

The implied diagnostic parameter $\hat{\theta} =
|\hat{\beta}_3|/\omega$ depends on both the scaling and the
assumed CPI weight $\omega$.
Under the baseline scaling and $\omega = 0.05$,
$\hat{\theta} \approx 121$.
Under the unit ordinal scaling, $\hat{\theta} \approx 126$;
under the wide scaling, $\hat{\theta} \approx 112$;
under the asymmetric scaling, $\hat{\theta} \approx 117$.

These values are all large relative to estimates of $\theta$
in the professional-forecast literature, where
\citet{Bordalo2020b} report values on the order of 0.5 to 1.5.
The inflation of $\hat{\theta}$ in this setting is expected:
the ordinal expectation measure is a coarse proxy for the
continuous belief that the diagnostic model assumes.
The three-category variable compresses the rich heterogeneity
in household expectations into three bins, amplifying the
apparent diagnostic parameter.
A continuous expectation measure (e.g., a point forecast in
percentage terms) would produce a lower $\hat{\theta}$
because it would capture gradations within the ``rise''
category that the ordinal variable conflates.

The scaling exercise confirms that the ordinal nature of the
expectation measure introduces quantitative imprecision in
the implied $\theta$ but does not affect the qualitative
conclusion.
The estimated forecast error exceeds the rational benchmark
by a factor of at least 50 under every scaling considered, and
the sign of $\hat{\beta}_3$ is invariant to the choice of
mapping.
